% Ovo je glavni fajl za generiranje predloska dokumenta doktorske disertacije Elektrotehnickog fakulteta u Sarajevu. 
% Verzija: 24. februar 2016. godine 

% Autor: Emir Sokic , esokic@etf.unsa.ba

%%%%%%%%%%%%%%%%%%%%%%%%%%%%%%%%%%%%%%%%%%%%%%%%%%%%%%%%%%%%%%%%%%%%%%
%%%%%%%%%%%%%%%%%%%%%%%%% POSTAVKE %%%%%%%%%%%%%%%%%%%%%%%%%%%%%%%%%%%
%%%%%%%%%%%%%%%%%%%%%%%%% NE MIJENJATI %%%%%%%%%%%%%%%%%%%%%%%%%%%%%%%
%%%%%%%%%%%%%%%%%%%%%%%%%%%%%%%%%%%%%%%%%%%%%%%%%%%%%%%%%%%%%%%%%%%%%%

\documentclass[12pt,oneside, a4paper]{book}

%Ovo su paketi koji se koriste za kompajliranje dokumenta
%Vecina paketa je ukljucena po defaultu u texlive 2013 i novijim verzijama
\usepackage{etex}
\usepackage{xcolor}
\usepackage[pdftex]{graphicx}
\usepackage{rotating}
\usepackage{epsfig}
\usepackage{epstopdf}
\usepackage[T1]{fontenc}
\usepackage[utf8]{inputenc}
\usepackage{cmap}
\usepackage[bosnian]{babel}
\usepackage[unicode]{hyperref}
%\usepackage{mathptmx}
\usepackage{newtxmath}
\usepackage{amscd}
\usepackage{amssymb}
\usepackage{amsmath}
\usepackage{amsfonts}
\usepackage[left=2.5cm,right=2.5cm,top=2.5cm,bottom=2.5cm]{geometry}
\usepackage{setspace} 
\usepackage{hhline}
\usepackage{enumerate}
\usepackage{delarray}
\usepackage{array}  
\usepackage{tabularx} 
\usepackage{multirow}  
\usepackage[bf, font=small]{caption}
\usepackage[labelfont=small, font=small]{subcaption}
\usepackage{wasysym}
\usepackage{subeqnarray}
\usepackage{pdflscape} % setting page into landscape view
\usepackage{enumitem} % for itemize lists
\usepackage[toc,page]{appendix}
\newcommand{\HRule}{\rule{\linewidth}{0.5mm}}
\usepackage{makeidx}
\usepackage{nomencl}
\usepackage{listings}
\lstset{basicstyle=\ttfamily,breaklines=true}
\usepackage{courier}
% Podesavanje izgleda zaglavlja i podnozja strana
\usepackage{fancyhdr} 

\fancypagestyle{plain}{%
  \fancyhf{}% Clear header/footer
  \fancyfoot[OR]{{\thepage}}%
  \fancyfoot[EL]{{\thepage}}%
  \renewcommand{\headrulewidth}{0pt}%
}

% required for printing index
% use \index{name} in text
%\usepackage{makeidx}
%\makeindex
% required for printing nomenclature
% use \nomenclature{symbol}{description} in text
%\usepackage{nomencl}
%\makenomenclature
%\renewcommand{\nomname}{Popis oznaka}




%Opcionalno
%\linespread{1.3}
%\setlist{nolistsep}   % setting for itemize lists
%\renewcommand{\thefootnote}{\fnsymbol{footnote}}  % to get unnumbered footnotes

% Adding a dot after chapter number in TOC 
%\let\savenumberline\numberline
%\def\numberline#1{\savenumberline{#1.}}



%\pagestyle{fancyplain}


%\rfoot{\thepage}
% iskljucivanje broja strane iz Sadrzaja, Popisa slika i Popisa tabela
\AtBeginDocument{\addtocontents{toc}{\protect\thispagestyle{empty}}}
\AtBeginDocument{\addtocontents{lof}{\protect\thispagestyle{empty}}}
\AtBeginDocument{\addtocontents{lot}{\protect\thispagestyle{empty}}}

%\rhead{\slshape \nouppercase \leftmark}
%\lhead{} %delete left header


%Podesavanje izgleda referenci
\usepackage[square, numbers, sort]{natbib} 

%Promjena naziva pojedinih poglavlja sa Hrvatskog na Bosanski
% Bibliography u "Literatura"
\addto\captionscroatian{%
  \renewcommand{\bibname}{Literatura}
  \renewcommand{\tablename}{Tabela}
  \renewcommand{\nomname}{Popis oznaka}
  \renewcommand{\indexname}{Indeks pojmova}
  \renewcommand{\lstlistingname}{Program}
}
%"Popis tablica" u "Popis tabela"
\addto\captionscroatian{\renewcommand{\listtablename}{Popis tabela}}
\addto\captionscroatian{\renewcommand\appendixname{Prilog}}
\addto\captionscroatian{\renewcommand\appendixpagename{Prilozi}}
\renewcommand\appendixtocname{Prilozi}

\makeindex
\makenomenclature

%\usepackage{etoolbox}
%\patchcmd{\chapter}{\thispagestyle{plain}}{\thispagestyle{fancyplain}}{}{}


\begin{document}

%%%%%%%%%%%%%%%%%%%%%%%%%%%%%%%%%%%%%%%%%%%%%%%%%%%%%%%%%%%%%%%%%%%%%%
%%%%%%%%%%%%%%%%%%%%%%%%% OSNOVNI DOKUMENT %%%%%%%%%%%%%%%%%%%%%%%%%%%
%%%%%%%%%%%%%%%%%%%%%%%%%%%%%%%%%%%%%%%%%%%%%%%%%%%%%%%%%%%%%%%%%%%%%%
%%%%%%%%%%%%%%%%%%%%%%%%%%%%%%%%%%%%%%%%%%%%%%%%%%%%%%%%%%%%%%%%%%%%%%

\frontmatter

%%%%%%%%%%%%%%%%%%%% NASLOVNA STRANA %%%%%%%%%%%%%%%%%%%%%%%%
\begin{titlepage}
\begin{center}

\includegraphics[width=0.25\textwidth]{etf-logo.png}~\\[0.1cm]
\textsc{\Large Univerzitet u Sarajevu}\\[0.2cm]  
\textsc{\Large Elektrotehnički fakultet}\\[0.2cm] 
\textsc{\Large Odsjek za automatiku i elektroniku}\\[3cm]\HRule \\[0.5cm] 
{\huge \bfseries Predložak za izradu doktorske disertacije - uz korištenje Latexa kao alata} \\[0.4cm] 
\HRule \\[0.5cm]

\textsc{\Large Doktorska disertacija}\\[0.4cm]
\textsc{\Large - Treći ciklus studija - }\\[1.5cm]

% Author and supervisor 
\textbf{ 
\Large Autor:\\  
\Large mr.sci. Nikola Tesla, dipl.ing.el.\\[1cm]  
\Large Mentor: \\[0.2cm] 
\Large Vanr.prof.dr Alessandro Volta, dipl.ing.el.} 
\vfill

% Bottom of the page  
{\large Sarajevo, \\februar 2016.}

\end{center} 
\end{titlepage}

%%%%%%%%%%%%%%%%%%%% POSTAVKA RADA %%%%%%%%%%%%%%%%%%%%%%%%%%%%%%%
\thispagestyle{plain}
\begin{flushleft}
\textbf{Vanr. prof. dr Alessandro Volta, dipl.el.ing}\\
\textbf{Viši ass. mr Thomas Edison, dipl.el.ing}\\
\textbf{Odsjek za automatiku i elektroniku}\\
\textbf{Sarajevo, 20.01.2015.}\\
\end{flushleft}

\begin{center}
\vspace{2cm}
{\Large Postavka zadatka završnog rada II ciklusa:}\\
\vspace{0.2cm}
{\large \textbf{Predložak za izradu doktorske disertacije - uz korištenje Latexa kao alata}}
\end{center}

\vspace{0.5cm}
U okviru rada je potrebno razviti \LaTeX\ predložak za izradu završnog rada prvog i drugog ciklusa studija, magistarskog rada po Ante-Bologna sistemu, te doktorske disertacije. U radu je potrebno:
\begin{itemize}
\item objasniti opći postupak za izradu i pisanje završnih radova prvog i drugog ciklusa, magistarskih radova po Ante-Bologna principima studiranja, te doktorskih disertacija poštivajući odgovarajuće Pravilnike,
\item dati kratko uputstvo za korištenje .tex dokumenata, pisanje slika, tabela i relacija, generisanje sadržaja, popisa slika i tabela i sl.
\item dati kratko upusto za generiranje odgovarajućih .pdf dokumenata iz odgovarajućih .tex fajlova,
\end{itemize}
Preporučuje se korištenje TexLive \LaTeX\ podrške, verzije 2013 ili novije.


\vspace{1cm}
\textbf{Polazna literatura:}

\begin{itemize}
\item[]  [1] Sonka, M., Hlavec, V., Boyle, R., "Image Processing, Analysis, and Machine Vision", Thomson, 2008.
\item[] [2] Lee, E.A., Varaiya, P., "Structure and Interpretation of Signals and Systems", Electrical Engineering \& Computer Science University of California, Berkeley, 2000

\end{itemize}


\begin{center}
\vspace{2cm}
\makebox[8cm]{\hrulefill} \\
Vanr. prof. dr Alessandro Volta, dipl. ing. el.\\
\end{center}





%%%%%%%%%%%%%%%%%%%% POTPISI CLANOVA KOMISIJE %%%%%%%%%%%%%%%%%%%%
\pagebreak
\thispagestyle{plain}
\vspace{2cm}
\begin{center}
\textbf{\Large Potpisi članova Komisije za ocjenu i odbranu doktorske disertacije}\\


\vspace{4cm}
\makebox[10cm]{\hrulefill} \\
Red. prof. dr Albert Einstein, dipl.el.ing, \\
predsjednik komisije\\

\vspace{3cm}
\makebox[10cm]{\hrulefill} \\
Van. prof. dr Alessandro Volta, dipl.el.ing, \\
mentor i član komisije\\

\vspace{3cm}
\makebox[10cm]{\hrulefill} \\
Doc. dr Michael Faraday, dipl.el.ing, \\
član komisije

\end{center}

%%%%%%%%%%%%%%%%%%%% PREDGOVOR %%%%%%%%%%%%%%%%%%%%

\chapter*{Predgovor}

U predgovoru autor ima pravo da izrazi svoja lična uvjerenja i dojmove koji se vežu za sam rad. Najčešće se u predgovoru pišu zahvalnice svima koji su direktno ili indirektno pomogli prilikom pisanja rada (porodica, profesori i kolege itd.). 

Nadam se da će ovaj predložak za izradu disertacije pomoći svim studentima Elektrotehničkog fakulteta u Sarajevu da napišu što bolje i kvalitetnije završne radove prvog, drugog ili trećeg ciklusa.

\vspace{1cm}
\begin{flushright}
\begin{tabular}{m{6cm}}
Nikola Tesla,\\
februar 2016.
\end{tabular}
\end{flushright}


%%%%%%%%%%%%%%%%%%%% IZJAVA AUTORA %%%%%%%%%%%%%%%%%%%%

\chapter*{Izjava autora}

Ja dolje potpisani izjavljujem da ovaj podnesak predstavlja moj vlastiti rad i da, prema mom saznanju i vjerovanju, u njemu nije sadržan bilo koji dio koji je prethodno objavljen ili napisan od strane drugih osoba, niti materijal koji je prihvaćen u svrhu dodjele diplome bilo kojeg stupnja na bilo kojoj visokoškolskoj ili naučnoj instituciji u svijetu.\footnote{Ova izjava je napisana u skladu sa Pravilnikom o strukturi i sadržaju doktorske disertacije i magistarskog rada na Elektrotehničkom fakultetu u Sarajevu (br. 04-1-673/11, dana 17.01.2011. godine).

Ipak, obzirom na ozbiljnost izrade jednog naučnog rada, autor zadržava pravo da radi komparacije koristi i izvore koji nisu njegovo originalno djelo, ali su obavezno navedeni u tekstu i referencama rada.}

Autorska prava doktorske disertacije ostaju pri autoru disertacije i Elektrotehničkom fakultetu Univerziteta u Sarajevu.

\begin{center}
\vspace{3cm}
\makebox[8cm]{\hrulefill} \\
mr. sci. Nikola Tesla, dipl. ing. el.\\
\vspace{8cm}
U Sarajevu,\\
februar 2016.
\end{center}

%%%%%%%%%%%%%% SAŽETAK I ABSTRACT %%%%%%%%%%%%%%%%%%%%%%%%%%%%%%%%%%%%%%%%%
%Sažetak (na Bosanskom jeziku) i Abstract (na Engleskom jeziku) su dati kao odvojeni fajlovi

\include{abstrakti/abstract_bh}
\include{abstrakti/abstract_en}

%%%%%%%%%%%%%%% SADRŽAJ %%%%%%%%%%%%%%%%%%%%%%%%%%%%%%%%%%%%%%%%%%%%%%%%%%%
%\clearpage
\tableofcontents
%%%%%%%%%%%%%%% POPIS SLIKA %%%%%%%%%%%%%%%%%%%%%%%%%%%%%%%%%%%%%%%%%%%%%%%
%\clearpage
\listoffigures
\addcontentsline{toc}{chapter}{Popis slika}
%%%%%%%%%%%%%%% POPIS TABELA %%%%%%%%%%%%%%%%%%%%%%%%%%%%%%%%%%%%%%%%%%%%%%
%\clearpage
\listoftables
\addcontentsline{toc}{chapter}{Popis tabela}
%%%%%%%%%%%%%%%%%%%%%%%%%%% INDEKS POJMOVA %%%%%%%%%%%%%%%%%%%%%%%%%%%%%%%%
\addcontentsline{toc}{chapter}{Indeks pojmova}
\printindex 
%%%%%%%%%%%%%%% POPIS OZNAKA%%%%%%%%%%%%%%%%%%%%%%%%%%%%%%%%%%%%%%%%%%%%%%%
%\clearpage
%\printnomenclature
%\addcontentsline{toc}{chapter}{Popis oznaka}


\cleardoublepage % start new page
\pagestyle{fancyplain} % puts headers/footers back on
\fancyhf{}
\lhead{\nouppercase{\fancyplain{}{\leftmark}}}
\renewcommand{\chaptermark}[1]{\markboth{#1}{}}
\renewcommand{\footrulewidth}{0.4pt} %draw foot line
\lfoot{\slshape Tesla, N., "Predložak za izradu doktorske disertacije"}
\rfoot{\thepage}
\cfoot{}

%%%%%%%%%%%%%%%%%%%%%%%%%%%%%%%%%%%%%%%%%%%%%%%%%%%%%%%%%%%%%%%%%%%%%%%%%%%
\mainmatter


%%%%%%%%%%%%%%%%%%%%%%%%%%%%%%%%%%%%%%%%%%%%%%%%%%%%%%%%%%%%%%%%%%%%%%
%%%%%%%%%%%%%%%%%%%%%%%%% POGLAVLJA %%%%%%%%%%%%%%%%%%%%%%%%%%%%%%%%%%
%%%%%%%%%%%%%%%%%%%%%%%%%%%%%%%%%%%%%%%%%%%%%%%%%%%%%%%%%%%%%%%%%%%%%%
%%%%%%%%%%%%%%%%%%%%%%%%%%%%%%%%%%%%%%%%%%%%%%%%%%%%%%%%%%%%%%%%%%%%%%

%Poglavlja je najbolje raditi u odvojenim fajlovima
%Poglavlje 1
\include{poglavlja/poglavlje_1}
%Poglavlje 2
\include{poglavlja/poglavlje_2}
%Poglavlje 3
\include{poglavlja/poglavlje_3}


%%%%%%%%%%%%%%%%%%%%%%%%%%%%%%%%%%%%%%%%%%%%%%%%%%%%%%%%%%%%%%%%%%%%%%
%%%%%%%%%%%%%%%%%%%%%%%%% PRILOZI %%%%%%%%%%%%%%%%%%%%%%%%%%%%%%%%%%%%
%%%%%%%%%%%%%%%%%%%%%%%%%%%%%%%%%%%%%%%%%%%%%%%%%%%%%%%%%%%%%%%%%%%%%%
%%%%%%%%%%%%%%%%%%%%%%%%%%%%%%%%%%%%%%%%%%%%%%%%%%%%%%%%%%%%%%%%%%%%%%
\begin{appendices}
%Priloge je najbolje raditi u odvojenim fajlovima
%Prilog 1
\include{prilog/prilog_1}
\end{appendices}




%%%%%%%%%%%%%%%%%%%%%%%%%%%%%%%%%%%%%%%%%%%%%%%%%%%%%%%%%%%%%%%%%%%%%%%%%%%
\backmatter
%%%%%%%%%%%%%%%%%%%% LISTA PUBLIKACIJA %%%%%%%%%%%%%%%%%%%%%%%%%%%%%%%%%%%%
\chapter*{Lista publikacija}
\addcontentsline{toc}{chapter}{Lista publikacija}

%%%%%%%%%%%%%%%%%%%%%%% LITERATURA %%%%%%%%%%%%%%%%%%%%%%%%%%%%%%%%%%%%%%%%
\addcontentsline{toc}{chapter}{Literatura}
\bibliographystyle{IEEEtranETF} 
\bibliography{literatura}



\end{document}
